\section{Kanban, JIT y Defectos}

\begin{ejercicio}{Guia 10, pregunta 2a y 2d}
La empresa ABC produce tarjetas telefonicas con demanda de $200$ tarjetas por hora. El proceso tiene tres operaciones:
\[
\begin{array}{c|cccc}
\text{Operacion} & C_i & t_{p,i} & t_{k,i} & t_{v,i}\\
\hline
P1&200&85&45&200\\
P2&250&78&67&300\\
P3&300&50&92&150
\end{array}
\]
donde los tiempos estan en segundos y $C_i$ es el tamano del contenedor/lote. El proceso descarta inicialmente $15\%$ de las unidades. Luego se pregunta que pasa si el descarte baja a $5\%$.
\end{ejercicio}

\begin{materia}{numero de kanbans}
Para cada operacion $i\in\{1,2,3\}$:
\[
  L_i=t_{p,i}+t_{k,i}+t_{v,i},
\]
donde $L_i$ es el tiempo total de reposicion asociado a la operacion $i$.

Sea:
\begin{itemize}
  \item $D$ la demanda de unidades buenas por unidad de tiempo.
  \item $y$ el rendimiento del proceso, es decir, la fraccion de unidades producidas que no se descarta.
  \item $C_i$ el numero de unidades por contenedor de la operacion $i$.
  \item $N_i$ el numero de kanbans requeridos en la operacion $i$.
\end{itemize}
Si la demanda $D$ esta expresada como demanda de unidades buenas, la produccion bruta requerida es $D/y$. Por eso:
\[
  N_i=\frac{D L_i}{y C_i}.
\]
La idea economica es directa: si baja el rendimiento, se necesitan mas unidades en circulacion para sostener la misma salida buena.
\end{materia}

\begin{solucion}{calculo base}
Los tiempos de reposicion son:
\[
  L_1=85+45+200=330,\qquad
  L_2=78+67+300=445,\qquad
  L_3=50+92+150=292.
\]
Con descarte de $15\%$, el rendimiento es:
\[
  y=1-0.15=0.85.
\]
La guia reporta los siguientes kanbans:
\[
  N_1=389,\qquad N_2=419,\qquad N_3=230.
\]
Lo importante para examen es la estructura:
\[
  N_i \propto \frac{L_i}{yC_i}.
\]
Si $L_i$ sube, se necesitan mas kanbans. Si $C_i$ sube, cada kanban mueve mas unidades y se necesitan menos. Si el rendimiento $y$ sube, se necesitan menos kanbans.
\end{solucion}

\begin{solucion}{si el descarte baja a $5\%$}
Ahora:
\[
  y_{\text{nuevo}}=1-0.05=0.95.
\]
Como todo lo demas queda constante:
\[
  N_{i,\text{nuevo}}
  =
  N_{i,\text{viejo}}\frac{0.85}{0.95}.
\]
Reemplazando los valores de la pauta:
\[
\begin{aligned}
N_{1,\text{nuevo}}&=389\frac{0.85}{0.95}\approx 348,\\
N_{2,\text{nuevo}}&=419\frac{0.85}{0.95}\approx 375,\\
N_{3,\text{nuevo}}&=230\frac{0.85}{0.95}\approx 206.
\end{aligned}
\]
La conclusion no es solo numerica: mejorar calidad reduce inventario en proceso, reduce tarjetas kanban necesarias y hace mas liviano el sistema JIT.
\end{solucion}

\begin{alerta}{unidades}
La formula exige unidades consistentes. Si $D$ esta en unidades/hora, entonces $L_i$ debe estar en horas. Si $L_i$ esta en segundos, convierte $D$ a unidades/segundo o $L_i$ a horas. En la guia OCR los resultados estan dados directamente; para examen, cuida esta conversion porque es donde se pierden puntos faciles.
\end{alerta}
